\documentclass{beamer}

\usepackage[spanish]{babel}

\usepackage{amsmath}
\usepackage{color}
\usepackage{minted}
\usepackage{hyperref}
\usepackage{multicol}
\usepackage{tabularx}
\usepackage{tikz}

\usetikzlibrary{shapes,arrows,positioning,shadows}

\newcommand{\bftt}[1]{\textbf{\texttt{#1}}}
\newcommand{\comment}[1]{{\color[HTML]{008080}\textit{\textbf{\texttt{#1}}}}}
\newcommand{\cmd}[1]{{\color[HTML]{008000}\bftt{#1}}}
\newcommand{\bs}{\char`\\}
\newcommand{\cmdbs}[1]{\cmd{\bs#1}}
\newcommand{\lcb}{\char '173}
\newcommand{\rcb}{\char '175}
\newcommand{\cmdbegin}[1]{\cmdbs{begin\lcb}\bftt{#1}\cmd{\rcb}}
\newcommand{\cmdend}[1]{\cmdbs{end\lcb}\bftt{#1}\cmd{\rcb}}

\newenvironment{exampletwouptiny}
  {\VerbatimEnvironment
   \begin{VerbatimOut}{example.out}}
  {\end{VerbatimOut}
   \setlength{\parindent}{0pt}
   \fbox{\begin{tabular}{l|l}
   \begin{minipage}{0.55\linewidth}
     \inputminted[fontsize=\scriptsize,resetmargins]{latex}{example.out}
   \end{minipage} &
   \begin{minipage}{0.35\linewidth}
     \setlength{\parskip}{6pt plus 1pt minus 1pt}%
     \raggedright\scriptsize\input{example.out}
   \end{minipage}
   \end{tabular}}}

\title{\LaTeX}
\author{DMCC poner los logos}

\begin{document}
\begin{frame}
    \maketitle    
\end{frame}

\begin{frame}{Qué es}
    Un sistema de preparación de documentos.
    Es un  sistema de composición tipográfica de alta calidad. Diseñado para la producción de documentación técnica y científica. \LaTeX{} está disponible como software libre.
\end{frame}


\begin{frame}{Porqué \LaTeX}
   \begin{itemize}
       \item Produce documentos elegantes.
       \item Está diseñado por y para científicos con una activa comunidad que lo mantiene. 
       \item Es poderoso y versátil.
   \end{itemize}
\end{frame}

\begin{frame}[fragile]{¿Cómo funciona?}
\begin{itemize}
\item Escribir el documento en texto plano, con comandos que le dan la estructura y significado. 
\item Guardar el documento.
\item EL programa de \LaTeX{} toma ese documento, lo procesa y produce un documento final con el formato y estructura que le pedimos.
\end{itemize}
\vskip 2ex
\begin{center}
\begin{minted}[frame=single]{latex}
Esto es \emph{texto plano} y esto una variable $x$ 
\end{minted}
\vskip 2ex
\tikz\node[single arrow,fill=gray,font=\ttfamily\bfseries,%
  rotate=270,xshift=-1em]{latex};
\vskip 2ex
\fbox{Esto es \emph{texto plano} y esto una variable $x$}
\end{center}
\end{frame}

\begin{frame}{Como usarlo}
    Uno puede usar \LaTeX:
    \begin{itemize}
        \item Local: Descargando una distribución de \LaTeX, y escribiendo en el editor de texto favorito. \item Online: En nuestro caso usaremos Overleaf. 
    \end{itemize}
\end{frame}

\begin{frame}[fragile]{Estructura de un documento}
   
    \begin{minted}[frame=single]{latex}
    \documentclass{article}
    \usepackage[spanish]{babel}
    
    \begin{document}
    El contenido %cometario
    \end{document}
    \end{minted}
\end{frame}

\begin{frame}{El lenguaje de \LaTeX}
    \begin{itemize}
        \item Los comandos parten con \cmdbs{}
        \item Todos los documentos deben partir con el comando \cmdbs{documentclass}
        \item Los argumentos de un comando se entregan entre \cmd{\{ }y \cmd{\} } 
        \item Algunos cáracteres tienen significados especiales en \LaTeX{} como \cmd{\$}, \cmd{\#}, \cmd{\%}, \cmd{\&} 
    \end{itemize}
\end{frame}

\begin{frame}[fragile]{Como escribir matemáticas}
    Para escribir matemáticas se debe usar \cmd{\$}

    \begin{exampletwouptiny}
Sean $a$ y $b \in \mathbb{N}$

$a_0 + a_1x + a_2x^2 + a_3x^3 + \ldots$
    \end{exampletwouptiny}
\vspace{.5cm}
    \begin{itemize}
        \item Es importante abrir y cerrar el ambiente matemático
        \item Existen comandos que solo funcionan dentro del ambiente matemático
    \end{itemize}
\end{frame}


\begin{frame}[fragile]{Como escribir matemáticas}
Hay veces que la ecuación es más grande y no corresponde ponerla en la misma linea que el texto, en estos casos podemos usar \cmdbs{[} y \cmdbs{]} para abrir y cerrar el entorno matemático:
      \begin{exampletwouptiny}
Esta es una sumatoria importante:
\[ S = \sum_{i=0}^n \frac{x^i}{i}\]
Con esto podemos concluir algo más 
importante ...
\end{exampletwouptiny}
\end{frame}

\begin{frame}[fragile]{Ambientes}
    \begin{itemize}
        \item Los ambientes son diferentes contextos de interpretación de comandos y formatos.
        \item Un mismo comando puede ser interpretado de manera distinta en distintos ambientes
        \item Todos los ambientes deben partir con \cmdbs{begin} y terminar con \cmdbs{end}
 \begin{exampletwouptiny}
\begin{enumerate}
    \item el item uno
    \item el item dos
\end{enumerate}

este texto 
\begin{center}
    y este texto
\end{center}
\end{exampletwouptiny}
\item El ambiente matemático también se puede usar con \cmdbegin{equation} y \cmdend{equation}
\end{itemize}
\end{frame}

\begin{frame}[fragile]{Paquetes}
\begin{itemize}
    \item \LaTeX{} viene con muchos comandos y ambientes por defecto. 
    \item Un paquete es una libreria de comandos y ambientes que se pueden importar. 
    \item Para importar un paquete se debe declarar en el preambulo con \cmdbs{usepackage}.
\end{itemize}
 \begin{minted}[frame=single]{latex}
    \documentclass{article}
    \usepackage[spanish]{babel}
    \usepackage{amsmath}
    
    \begin{document}
    El contenido %cometario
    \end{document}
    \end{minted}

\end{frame}

\begin{frame}{Ejemplos}

Buscar ejemplos diversos de paquetes en diferentes areas, biologia quimica ciencias de la compu. etc.
    
\end{frame}

\end{document}